\section{Particular inverses and the literal quadratic metric}

\subsection{A finite Laurent particular solution}

Put \(f=r^2-1\), \(\lambda_2=(2s+1)(2s+2)\), and define
\[
 D_s=\partial_r+\frac{2s}{rf},\qquad
 \mathcal L_{s,\mu}=D_s^2+\frac{2r}{f}D_s+
                      \frac{\mu}{r^2f}-\frac{4s^2}{f^2}.
\]
The twist is essential: it differentiates the retained radial envelope. Set
\begin{align*}
 B&=\frac{s^2(s+1)^2(\mu-2s^2-6s-6)}{8(2-\mu)},\\
 M_s(\mu)&=-8s+2s\mu+s\mu^2-\tfrac52s^2-\tfrac94s^2\mu
 +\tfrac52s^2\mu^2+\tfrac{79}2s^3-\tfrac{17}2s^3\mu\\
 &\hspace{8mm}+54s^4-\tfrac{29}4s^4\mu+\tfrac{51}2s^5+\tfrac{11}2s^6,\\
 A&=\frac{-M_s(\mu)/(2-\mu)-6B}{\mu-\lambda_2}.
\end{align*}
Here \(M_s(\mu)\) is a polynomial, unrelated to the finite-part source moment \(M(w,s,a)\).

\begin{theorem}[Scalar--scalar Laurent inverse]\label{thm:laurent}
If \(\mu\ne2,\lambda_2\), then \(Y=r+A/r+B/r^3\) satisfies the original reduced scalar--scalar source equation
\[
 \mathcal L_{s,\mu}Y+\mathcal Q_{s,\mu}=0\quad(r>1),
\]
where the source has the exact decomposition
\[
 r^2f\mathcal Q_{s,\mu}=-2r^3-(\mu+2s-4s^2)r
 +\frac{M_s(\mu)}{(2-\mu)r}
 -\frac{[\mu-(2s+3)(2s+4)]B}{r^3}.
\]
At either physical parent rate, the associated sourced reconstruction data exist for every independent angular value \(\mu=-k^2\), \(k\in\R\).
\end{theorem}
\cert{B1}
\begin{proof}
Twice differentiating the Laurent polynomial, including the twist, gives
\[
 r^2f\mathcal L_{s,\mu}Y=2r^3+(\mu+2s-4s^2)r
 +\frac{(\mu-\lambda_2)A+6B}{r}
 +\frac{[\mu-(2s+3)(2s+4)]B}{r^3}.
\]
The definition of \(A\) cancels the middle source coefficient; the other coefficients cancel directly. The source decomposition itself follows by clearing denominators in the original source rows, with no removal of their trace-free term. The Laurent function is smooth on \(r>1\). For physical rates \(\re s<-1\), \(\im s\ne0\), and
\(\im\lambda_2=(8\re s+6)\im s\ne0\). Thus \(-k^2\ne\lambda_2\), and \(-k^2\ne2\); also \(s\ne0\). These are exactly the denominator and regularity requirements of the sourced reconstruction.
\end{proof}

At the fundamental harmonic \(\mu=(s+1)(s+2)=\rho(\rho-1)\), the first inverse-power coefficients simplify to
\begin{align*}
 A^{SS}_+(\rho)&=-\frac{(5\rho-7)(\rho+2)}{2(\rho+1)},&
 A^{SS}_-(\rho)&=\frac{(5\rho+2)(\rho-3)}{2(\rho-2)},\\
 A^{VV}_+(\rho)&=\tfrac23(\rho-2)(4\rho+3),&
 A^{VV}_-(\rho)&=\tfrac23(\rho+1)(4\rho-7).
\end{align*}
The last row concerns the separate vector--vector particular field.
\cert{B2}
These formulas determine particular coefficients. They do not, by themselves, prove retarded uniqueness or admission of the forced growing term \(r\) in the desired boundary class. A previous claimed obstruction to rational scalar--scalar reconstruction was retracted in the working notes after correction of untwisted derivatives and a doubled trace-free coefficient. The positive survivor is the explicit inverse above.

\subsection{The nonseparated scalar--scalar metric}

The independent-harmonic construction is not the endpoint of the source. Let
\[
 F=S_\rho^2,\qquad Q=\code{coordinateResponse}(q),\qquad
 R=\frac{2}{\lambda_2}\code{fullAngularField}(q,b).
\]
These are constructed smooth angular fields, and their equations are
\begin{equation}\label{eq:threeangular}
 \Delta Q=2Q-F,\qquad \Delta R=\lambda_2R+F.
\end{equation}
In particular, they are not arbitrary inverse-oracle arguments to the final theorem. Appendix~\ref{app:metric} gives the radial coefficient recipe. For \(Y_0=F,Y_2=Q,Y_D=R\), form
\begin{equation}\label{eq:metric}
 K=a^2e^{2s\tau}(1-r^{-2})^s
 \sum_{k\in\{0,2,D\}}Y_k
 \begin{pmatrix}
 F_{0k}&F_{1k}&0&0\\
 F_{1k}&F_{2k}&0&0\\
 0&0&2r^2H_k/y^2&0\\
 0&0&0&2r^2H_k/y^2
 \end{pmatrix}.
\end{equation}

\begin{theorem}[Literal sourced Einstein equation]\label{thm:metric}
For every Xi zero copy, either branch, every complex amplitude, and every exterior point \(r>1,y>0\), the twice differentiable field \eqref{eq:metric} satisfies all sixteen coordinate equations
\[
 D\mathcal E_{g_0}[K]+\mathcal Q^{SS}_{q,b,a}=0,
 \qquad \mathcal E(g)=G(g)-3g.
\]
Here \(\mathcal Q^{SS}_{q,b,a}\) is the actual scalar--scalar quadratic Einstein source, normalized as one half of the second variation of the original first-order scalar metric.
\end{theorem}
\cert{B3}
\begin{proof}[Proof structure]
Differentiate each of the three actual metric columns, retaining its envelope. The coordinate linearized operator is linear in these columns. Substitute \eqref{eq:threeangular} and the scalar equation for \(S_\rho\). Its angular terms reduce to \(F\), \(\Delta F\), its first derivatives, and the trace-free Hessians of \(F\) and \(Q\). The source has exactly the same tensor structures. The radial recipe in Appendix~\ref{app:metric} makes each coefficient of their sum zero. The corpus proves this coefficient cancellation before substituting the actual smooth fields, then identifies their two-jets with the original coordinate Einstein source. Its final theorem quantifies over both coordinate indices; no sector equation is left as an argument. The full component algebra and angular inverse constructions are in the accompanying dependency source snapshot.
\end{proof}

This closes the local, nonseparated scalar--scalar equation at the stated source rate. It does not construct a switched, retarded quadratic metric in every sector, the opposite-branch correction, or its complete integrated cubic action. Consequently the remaining action problem cannot be described as merely finding this scalar--scalar particular metric, but it also cannot be declared solved by its existence.

\section{The cubic coefficient retains the changing test weight}

Write \(L,B,C\) for the first, second, and third derivatives of \(\mathcal E\) at \(g_0\). Consider the coefficient of \(\epsilon^2\eta\) in a metric family with terms
\[
 g_0+\epsilon p+\eta m+\epsilon^2k_{++}
       +\epsilon\eta k_{+-}+\epsilon^2\eta\ell.
\]
Let \(W(g)\) be the metric-dependent contraction against a fixed test tensor, \(W_0=W(g_0)\), and \(W_1[h]=DW_{g_0}[h]\). For the original linearized solutions \(p,m\), the exact local paired coefficient is
\begin{align}\label{eq:cubic}
 &W_0\!:\!\left(L\ell+B(p,k_{+-})+B(m,k_{++})+\tfrac12C(p,p,m)\right)\nonumber\\
 &\quad+W_1[m]\!:\!\left(Lk_{++}+\tfrac12B(p,p)\right)
       +W_1[p]\!:\!\left(Lk_{+-}+B(p,m)\right).
\end{align}
The colon denotes the full coordinate contraction. The source theorem requires the displayed correction fields to be twice differentiable at the point. Their nonlinear field equations are not required for \eqref{eq:cubic}.
\cert{D2}
Taylor expansion proves the formula: the cubic tensor term has multiplicity three in the symmetric third derivative, hence coefficient \(3/3!=1/2\); the two products of a first weight variation and a quadratic residual give the last line. Terms involving a first-order residual vanish because the original first-order modes solve the linearized equation.

\textbf{Conditional simplification.} If both \(Lk_{++}+B(p,p)/2=0\) and \(Lk_{+-}+B(p,m)=0\) hold, the last line vanishes. This is the explicit premise of \code{mixedPairedCoefficient_of_quadratic_zero}, not an unconditional physical extinction theorem.
\cert{D3}
For the affine family \(k_{++}=k_{+-}=\ell=0\), the corpus instead gives an unconditional formula retaining \(W_0:C/2\), \(W_1[m]:B(p,p)/2\), and \(W_1[p]:B(p,m)\). Dropping the latter two is incorrect. Passing from this local density to an integrated action further requires a measure, domain, and boundary prescription.

\section{Angular current and the renormalized modular boundary}

\subsection{A bilinear source-contact identity}

For the literal scalar field, abbreviate
\begin{align*}
 E&=y^2(S_x^2+S_y^2),&
 C_0&=\frac{3\lambda-2}{16}S^2-\frac E4,\\
 B_0&=y^2(S_x(C_0)_x+S_y(C_0)_y),&
 G_0&=y^2((C_0)_x^2+(C_0)_y^2).
\end{align*}
The full source-contact calculation yields
\begin{align}\label{eq:contact}
 &(\lambda-2)(\lambda-6)(3\lambda-2)S^4\nonumber\\
 &=16(\lambda-2)(\lambda-6)S^2C_0
 +64(\lambda-2)(2SB_0+C_0E)\nonumber\\
 &\hspace{8mm}-128E\Delta C_0-512G_0.
\end{align}
Its prefactor is nonzero at every Xi zero: \(\re\lambda=\sigma(\sigma-1)-t^2<0\), whereas its three possible roots are positive real numbers. The identity follows by applying the angular equation and product rule to the displayed contact and retaining all gradient terms.
\cert{D4}
All products in \eqref{eq:contact} are complex bilinear. Thus \(E\) and \(G_0\) need not be nonnegative. This rigidity identity cannot supply a positive integrated energy without further structure.

\subsection{The curved modular domain}

For \(Y\ge1\), use the actual truncated modular domain
\[
 \mathcal D_Y=\{(x,y):|x|\le1/2,\ x^2+y^2\ge1,\ y\le Y,\ y>0\}.
\]
At fixed null parameters \(p,v\), let \(F\) be the original scalar potential and let \(T=S_u\) be the testing Eisenstein scalar, \(\Delta T=\mu T\), \(\mu=u^2-1/4\). Define a coordinate flux
\[
 j=(\Delta F)\nabla T-\mu T\nabla F.
\]
Its Euclidean divergence is the source's coexact-vector defect. The measure in the next formula is \(\dd x\dd y\), not \(y^{-2}\dd x\dd y\).

\begin{theorem}[Original modular Green finite part]\label{thm:modular}
For actual source and reader Xi copies, either source branch and any amplitude,
\[
 \int_{\mathcal D_Y}\operatorname{div}j\dd x\dd y
       =\int_{-1/2}^{1/2}j_y(x,Y)\dd x.
\]
For fixed \(p,v\ge0\), subtract from this top flux the corpus's grouped nonconstant cusp-exponent family. Its limit at \(Y=e^k\), \(k\to\infty\), equals the coefficient of exponent zero. That coefficient, and hence this prescribed finite part, is zero when reader copy equals source copy or its functional-equation reflection.
\end{theorem}
\cert{B4}
\begin{proof}[Proof structure]
The scalar and coexact-vector identities identify the integrand with the coordinate divergence. Modular translation cancels the vertical sides. Modular inversion identifies the two pieces of the curved lower boundary with opposite flux; replacing the domain by a rectangle would lose this step. Stokes' formula leaves only the top edge. The actual Fourier--Bessel tail has the decay needed for its contribution to the subtracted limit to vanish. Grouping equal exponents before subtraction leaves the coefficient of exponent zero, including coincident-exponent cases. The diagonal and reflected zero-coefficient identities then give the final assertion. The formal statements apply to the original source, including onset and both Weyl terms.
\end{proof}

The finite part is a specified subtraction scheme, not the assertion that the unrenormalized integral converges to zero. Nor does a zero angular finite part at each fixed pair of parameters license exchanging that limit with an improper time integral or horizon limit.

\subsection{The actual original metric current}

There is also an exact local identification for the original physical angular current. In ingoing coordinates \((p,R,x,y)\), put \(U=RZ\), where \(Z\) is the original ingoing potential, and let \(W\) be the original left scalar with its physical radial normalization. For angular index \(j\), define
\[
 \mathcal B_j(F,G)=(\Delta F)\partial_jG-(\Delta G)\partial_jF.
\]
Then the original metric current is
\begin{equation}\label{eq:originalcurrent}
 J^j=\frac1{R^2}\left[(R^2-1)\mathcal B_j(U_R,W_R)
                 +\mathcal B_j(U_R,W_p)+\mathcal B_j(U_p,W_R)\right].
\end{equation}
The source and reader copies may be independent; both original testing sides and both radial normalizations are included. The proven domain is \(p>0,1<R<3,y>0\).
\cert{B5}
The proof commutes the actual mixed derivatives, identifies the causal stream as the coexact derivative of \(RZ\), and substitutes in the four-component Einstein current. The source and left-field smoothness are supplied by their constructions. In particular, differentiating \(Z\) while discarding the prefactor \(R\) does not reproduce this current.

Equation~\eqref{eq:originalcurrent} is a more faithful starting point than a projected current with angular components set to zero. Its derivative-weighted Green transfers still require uniform parameter-dependent boundary estimates before Theorem~\ref{thm:modular} can be used inside the complete spacetime action. That transfer is an open problem here.
